% Compile with LuaLaTeX:
% lualatex when_every_parenthesis_matters_AMM_note_v4_en.tex
% lualatex when_every_parenthesis_matters_AMM_note_v4_en.tex
\documentclass[11pt]{article}

\usepackage[margin=0.92in]{geometry}
\usepackage{amsmath,amssymb,amsthm}
\usepackage{booktabs}
\usepackage[hidelinks]{hyperref}

\newtheorem{theorem}{Theorem}
\newtheorem{lemma}[theorem]{Lemma}
\newtheorem{proposition}[theorem]{Proposition}

\newcommand{\Q}{\mathbb{Q}}

\title{When Every Parenthesis Matters}
\author{Takayuki Kuriyama}
\date{}

\begin{document}
\maketitle

\begin{abstract}
Fix the variable order $x_1,\ldots,x_n$ and form fully parenthesized
expressions using $+$, $\times$, and $/$. An inner pair of parentheses is
called necessary if deleting that pair alone, and then applying the usual
precedence and left-associativity conventions, changes the represented
rational function. We classify the expressions in which every inner pair is
necessary. A three-row local rule yields a recursive structure and the
Fibonacci enumeration
\[
 r_2=3,\qquad r_n=2F_{n+1}\quad(n\ge3).
\]
\end{abstract}

\section{The question}

Parentheses can be visibly present and mathematically idle. For example,
\[
 x_1+(x_2/x_3)=x_1+x_2/x_3,
\]
whereas
\[
 x_1/(x_2+x_3)\ne x_1/x_2+x_3.
\]
How many fully parenthesized expressions have no idle inner pair at all?

Fix distinct indeterminates $x_1,\ldots,x_n$ in this order. A
\emph{fully parenthesized expression} is a plane binary tree whose leaves,
from left to right, are $x_1,\ldots,x_n$, in which every internal vertex has
exactly two children and is labeled by
\[
 +,\qquad \times,\qquad /.
\]
We compare expressions as elements of $\Q(x_1,\ldots,x_n)$. For example, the
fully parenthesized expression in our sense corresponding to
$x_1+x_2\times x_3$ is $\bigl(x_1+(x_2\times x_3)\bigr)$. The outermost pair
is not a candidate pair, since deleting it changes neither the parsing nor the
represented rational function. Thus the $n-2$ candidate pairs are precisely
those corresponding to edges joining two internal vertices.

Whenever a candidate pair is deleted, the resulting infix expression is
parsed in the usual way: $\times$ and $/$ bind more strongly than $+$, and
operations at the same precedence level associate to the left.

A pair is \emph{necessary} if deleting that pair alone, while retaining all
other pairs, changes the rational function represented by the expression. An
expression is \emph{rigid} if all its $n-2$ candidate pairs are necessary.

For $n=3$, the six rigid fully parenthesized expressions are
\[
\begin{gathered}
 \bigl((x_1+x_2)\times x_3\bigr),\qquad
 \bigl((x_1+x_2)/x_3\bigr),\qquad
 \bigl(x_1\times(x_2+x_3)\bigr),\\
 \bigl(x_1/(x_2+x_3)\bigr),\qquad
 \bigl(x_1/(x_2\times x_3)\bigr),\qquad
 \bigl(x_1/(x_2/x_3)\bigr).
\end{gathered}
\]
Thus $r_3=6$.

For $n=2$, the three fully parenthesized expressions are
$(x_1+x_2)$, $(x_1\times x_2)$, and $(x_1/x_2)$. None has a candidate pair,
so all three are rigid vacuously and $r_2=3$.

Our objects are fully parenthesized binary trees, not minimally parenthesized
infix strings. Fixing the tree makes the necessity of a prescribed pair an
edge-local property. We restrict attention to the three operations above;
adding subtraction leads to a different enumeration problem.

\section{The local rule}

Call an internal child \emph{additive} if the operation at its root is $+$.
The following table determines whether the pair surrounding an internal child
is necessary.

\begin{table}[ht]
\centering
\small
\begin{tabular}{ccc}
\toprule
parent & left internal child & right internal child \\
\midrule
$+$      & never necessary & never necessary \\
$\times$ & necessary exactly when additive & necessary exactly when additive \\
$/$      & necessary exactly when additive & always necessary \\
\bottomrule
\end{tabular}
\caption{When the parentheses around an internal child are necessary.}
\label{tab:local}
\end{table}

\begin{lemma}[Local rule]\label{lem:local}
Let $v$ be a nonroot internal vertex of a fully parenthesized expression, and
let $p$ be its parent. Whether the pair surrounding $v$ is necessary depends
only on whether $v$ is the left or right child of $p$ and on the operation
labels at $p$ and $v$, as specified in Table~\ref{tab:local}.
\end{lemma}

\begin{proof}
In the local expressions below, we suppress the outermost pair of
parentheses.

Let $A$ and $B$ be the two children of $v$, and let $C$ be the sibling
subtree of $v$. Thus the relevant local expression is either
\[
 (A\circ B)\star C
 \qquad\text{or}\qquad
 A\star(B\circ C),
\]
where $\circ$ labels $v$ and $\star$ labels its parent.

First consider the entries declared unnecessary in
Table~\ref{tab:local}. Under the stated precedence and
left-associativity conventions, deleting the displayed pair gives the same
rational function in each of the following cases:
\[
\begin{aligned}
 (A\circ B)+C&=A\circ B+C,
 &A+(B\circ C)&=A+B\circ C
 &&(\circ\in\{+,\times,/\}),\\
 (A\times B)\times C&=A\times B\times C,
 &(A/B)\times C&=A/B\times C,\\
 A\times(B\times C)&=A\times B\times C,
 &A\times(B/C)&=A\times B/C,\\
 (A\times B)/C&=A\times B/C,
 &(A/B)/C&=A/B/C.
\end{aligned}
\]

For the remaining entries, the following positive specializations show that
deleting the pair changes the value. The triple in the first column records
the parent operation, the side of the internal child, and the operation at
that child.
\[
\begin{array}{c|c|c}
\text{case} & \text{before deletion} & \text{after deletion}\\
\hline
(\times,L,+) &(1+1)\times2=4&1+1\times2=3\\
(\times,R,+) &2\times(1+1)=4&2\times1+1=3\\
(/,L,+) &(1+1)/2=1&1+1/2=3/2\\
(/,R,+) &1/(1+1)=1/2&1/1+1=2\\
(/,R,\times) &1/(1\times2)=1/2&1/1\times2=2\\
(/,R,/) &1/(1/2)=2&1/1/2=1/2
\end{array}
\]

These numerical witnesses apply to subexpressions of arbitrary size. Indeed,
for any nonempty subexpression and any prescribed positive rational value
$q$, its variables can be assigned positive rational values so that the
subexpression evaluates to $q$: at a $+$-node use child values $q/2,q/2$,
and at a $\times$- or $/$-node use $q,1$. Since $A,B,C$ have disjoint
variable sets, their values can be prescribed independently.

It remains only to check that a local change cannot disappear higher in the
tree. Put $K=\Q(x_1,\ldots,x_n)$. Every subexpression represents a nonzero
element of $K$, since assigning $1$ to every variable gives it a positive
value. At each step toward the root, the current rational function $u$ is
transformed, for some nonzero sibling function $h$, by one of
\[
 u\longmapsto u+h,\qquad
 u\longmapsto hu,\qquad
 u\longmapsto u/h,\qquad
 u\longmapsto h/u.
\]
Each map is injective on its domain. Hence two distinct local functions remain
distinct all the way to the root. This proves the table.
\end{proof}

\begin{proposition}[Context independence]\label{prop:context}
Let $T$ be a subtree of a fully parenthesized expression $E$, and let $\pi$ be
a candidate pair inside $T$. Then $\pi$ is necessary in $T$ if and only if it
is necessary in $E$.
\end{proposition}

\begin{proof}
If deleting $\pi$ does not change the rational function represented by $T$,
it does not change the rational function represented by $E$. Conversely, if
deleting $\pi$ changes the function represented by $T$, the injectivity
argument in the proof of Lemma~\ref{lem:local} shows that the difference
persists from the parent of $T$ to the root. Hence the function represented by
$E$ also changes.
\end{proof}

\section{Structure and the Fibonacci numbers}

\begin{proposition}[Structural classification]\label{prop:structure}
A fully parenthesized expression with at least two leaves is rigid if and only
if exactly one of the following holds.
\begin{enumerate}
\item Its root is $+$ and both children are leaves.
\item Its root is $\times$ and each child is either a leaf or a two-leaf sum.
\item Its root is $/$, its left child is either a leaf or a two-leaf sum, and
      its right child is either a leaf or a rigid expression.
\end{enumerate}

Here a \emph{leaf} means a vertex labeled by a variable $x_i$, with no
operation symbol, in the full binary tree representing the expression.
\end{proposition}

\begin{proof}
Suppose that the expression is rigid. If its root is $+$, the local rule
prohibits every internal child, so both children are leaves. If its root is
$\times$, every internal child must be additive. By
Proposition~\ref{prop:context}, such a child is itself rigid, and the first
case then shows that it is a two-leaf sum. If the root is $/$, the same
argument applies to the left child. The local rule places no restriction on
the root operation of an internal right child, but that right subtree is rigid
by Proposition~\ref{prop:context}.

Conversely, consider an expression of one of the three listed forms. Every
pair immediately below the root is necessary by the local rule. Every pair
inside a rigid child is necessary in the whole expression by
Proposition~\ref{prop:context}. Hence the whole expression is rigid.
\end{proof}

\begin{theorem}\label{thm:main}
Let $F_1=F_2=1$. Then
\[
 r_2=3,
 \qquad
 r_n=2F_{n+1}\quad(n\ge3).
\]
Thus
\[
 r_2,r_3,r_4,r_5,r_6,r_7,r_8,\ldots
   =3,6,10,16,26,42,68,\ldots .
\]
\end{theorem}

\begin{proof}
We have already established $r_2=3$ and $r_3=6$.

For four leaves, there is one multiplication-rooted expression, with a
two-leaf sum on each side. Division contributes
\[
 r_3+r_2=6+3=9
\]
expressions, according as the left block has one or two leaves. Hence
$r_4=10$.

Now let $n\ge5$. A $+$-rooted rigid expression has only two leaves, and a
$\times$-rooted rigid expression has at most four leaves. Hence the root must
be $/$. Its left child is either the leaf $x_1$ or the two-leaf sum
$x_1+x_2$, while its right child is an arbitrary rigid expression on the
remaining variables. Therefore
\[
 r_n=r_{n-1}+r_{n-2}\qquad(n\ge5).
\]
Since $r_3=6=2F_4$ and $r_4=10=2F_5$, the stated formula follows.
\end{proof}

\section*{Perspective}

Representative work on redundant parentheses concerns syntactic grammars and
unparsing under precedence and associativity
\cite{Ramsey,Bates,Lotfallah}. Related enumerative work studies equivalence
classes of arithmetic expressions
\cite{StosicDamnjanovicRandelovic,Cohen}. The present problem is different:
the underlying fully parenthesized tree is fixed, and every prescribed inner
pair must be semantically necessary.

In OEIS notation,
\[
 r_n=A118658(n+2)\qquad(n\ge3).
\]
Thus Theorem~\ref{thm:main} gives a parenthesis interpretation of this
sequence \cite{OEIS}. OEIS A155069 concerns a different model, counting
minimally parenthesized infix strings formed with $+$ and $-$.

\begin{thebibliography}{9}
\footnotesize
\setlength{\itemsep}{0pt}

\bibitem{Bates}
R. M. Bates,
Distinguishing redundant parentheses by syntax,
\emph{Internat. J. Comput. Math.} \textbf{82} (2005), 951--967.

\bibitem{Cohen}
B. Cohen,
The number of non-isomorphic arithmetic expressions that can be constructed
using $+,-,\times$, and $/$,
arXiv:2602.18522 (2026).

\bibitem{Lotfallah}
W. B. Lotfallah,
Characterizing unambiguous precedence systems in expressions without
superfluous parentheses,
\emph{Internat. J. Comput. Math.} \textbf{86} (2009), 1--20.

\bibitem{Ramsey}
N. Ramsey,
Unparsing expressions with prefix and postfix operators,
\emph{Software: Practice and Experience} \textbf{28} (1998), 1327--1356.

\bibitem{StosicDamnjanovicRandelovic}
I. Sto\v{s}i\'c, I. Damnjanovi\'c, and \v{Z}. Ran{\dj}elovi\'c,
Counting the number of inequivalent arithmetic expressions on $n$ variables,
\emph{Filomat} \textbf{39} (2025), 949--962.

\bibitem{OEIS}
N. J. A. Sloane and The OEIS Foundation Inc.,
The On-Line Encyclopedia of Integer Sequences, A118658 and A155069.

\end{thebibliography}

\end{document}
